\section{Experiments}
\label{sec:experiments}

\subsection{Experimental Setup}

\paragraph{Research questions and evaluation overview.}
We evaluate three complementary questions. \emph{RQ1} asks whether
relation-aware graph retrieval improves complete evidence-set recovery
on large-scale multi-hop benchmarks under a fixed top-$k$ budget.
\emph{RQ2} asks whether EPGM recovers complete output dependencies in
controlled execution-provenance trajectories, and which structural
components---graph message passing, relation types, calibrated edge
weights, or inference-time target
promotion---are associated with the gain. \emph{RQ3} asks, via qualitative cases on a real multi-agent trace,
whether the same schema exposes support and revision structure that flat
retrieval buries. Table~\ref{tab:eval-overview} summarizes
how the three settings trade off scale, control, and realism; they are
complementary sources of evidence; only RQ1--RQ2 use automatic ranking
tables as primary results.

\begin{table*}[t]
\centering
\small
\caption{Evaluation settings. The three protocols trade off scale,
experimental control, and fidelity to real multi-agent execution.}
\label{tab:eval-overview}
\begin{tabular}{lp{3.2cm}ccp{4.2cm}}
\hline
Setting & Graph realization & Scale & Native execution & Primary role \\
\hline
Multi-hop QA & Evidence nodes + four rule-based edge types & High & No &
Scalable complete-support evaluation \\
Synthetic provenance & Task/Agent/Call/Output + \texttt{feeds} & High & Simulated &
Controlled dependency learning and ablation \\
Real multi-agent trace & Full EPGM schema + audit edges & Low & Yes &
Qualitative case analysis \\
\hline
\end{tabular}
\end{table*}

\paragraph{Datasets and splits.}
The evidence-path experiments use HotpotQA, 2WikiMultiHopQA, and
MuSiQue-Ans v1.0. For each dataset, we retain its training set,
deterministically shuffle the labeled development set with split seed 13,
use the first 500 shuffled examples for model selection, and reserve the
remaining shuffled examples as a held-out evaluation split. We additionally
construct a controlled execution-provenance benchmark from
2WikiMultiHopQA. It maps each recoverable ordered pair of supporting facts
to typed tool calls, tool outputs, and one output dependency. The converter keeps at most 32 candidate outputs per task and uses a pinned
hybrid successor scorer with equal BM25/dense weight. Each source ranks
non-self candidates from the question plus source evidence. For every
source output, schema v3 creates exactly two ordinary \texttt{feeds}
successors: one rank-1 semantic head and one deterministic branch from the
near (ranks 2--4), mid (5--8), or tail (9+) bucket. The required dependency
is inserted as an ordinary head or branch, and a non-gold branch in the
same rank bucket is required when the gold edge is not the head; no
gold-only field enters the ranking input. Source-local confidence is
softmax-normalized with temperature 0.1 and mapped to $0.5+0.5p$, so the
two \texttt{feeds} weights sum to 1.5. This
label-conditioned benchmark places the required path in the inference graph
alongside matched distractors and tests whether retrieval ranks and decodes
it. It is not intended to measure path discovery from unlabeled logs or to
substitute for naturally occurring multi-agent traces. Table~\ref{tab:data} gives the split sizes.
Evidence-path models select checkpoints on the 500-example selection
split. Dense-FT selects its encoder by development MAP@100; the R-GCN then
selects its checkpoint by
$0.50\,\mathrm{FS@5}+0.30\,\mathrm{Recall@5}+0.20\,\mathrm{MRR}$. The synthetic
execution-provenance R-GCN uses the complete 2,850-item generated
development split for joint-score selection and the same 2,851-item
held-out test split for final reporting. For RQ3 we use SupplyChainTrace (one real multi-agent supply-chain
security session; internal id task2) with 150 retrievable nodes. We manually
annotate 20 exploratory queries but do not treat this single-session set as
a statistical benchmark. Instead, we inspect one support case and one
revision case to illustrate the two mechanisms; their selection is
illustrative rather than a significance test. No real-trace example is used
for training or parameter selection.

\begin{table*}[t]
\centering
\small
\caption{Datasets, split sizes, and primary evaluation targets.}
\label{tab:data}
\begin{tabular}{lrrrrl}
\hline
Dataset & Train & Selection & Evaluation & Top-$k$ & Primary target \\
\hline
HotpotQA & 90,025 & 500 & 6,869 & 10 & Evidence set and connectivity \\
2WikiMultiHopQA & 167,454 & 500 & 12,076 & 10 & Evidence set and dependency path \\
MuSiQue-Ans v1.0 & 19,938 & 500 & 1,917 & 10 & Evidence set and multi-hop path \\
Synthetic 2Wiki Provenance & 74,012 & 2,850 & 2,851 & 10 & Outputs and execution dependencies \\
SupplyChainTrace (1 session) & -- & -- & qualitative & -- & Case analysis only \\
\hline
\end{tabular}
\end{table*}

\paragraph{Compared methods.}
For evidence-path retrieval, \textbf{BM25} and \textbf{Dense} are
lexical and pretrained dense baselines. \textbf{BM25-Graph} and
\textbf{Dense-Graph} are legacy fixed heuristic rerankers retained for
comparison; they are distinct from the current FastGraphRAG-style baseline
used in RQ2. \textbf{Dense R-GCN} uses frozen dense
representations as node features in a relational graph convolutional
network. \textbf{Dense-FT} fine-tunes the dense retriever on the task,
whereas \textbf{Dense-FT R-GCN} initializes the R-GCN with the
fine-tuned representations and is our full model for this setting.

For controlled synthetic trajectories, \textbf{BM25}, \textbf{Dense}, and
\textbf{Dense-FT} provide sparse, pretrained dense, and task-adapted
dense baselines. \textbf{FastGraphRAG-style} is a compact, non-LLM
entity-bridge baseline: it deterministically extracts title/body entity
mentions, starts from dense seeds, resolves one candidate per title group,
and stably inserts accepted bridge targets. An accepted source-to-target
entity bridge is emitted as a directed \texttt{bridge\_to} edge only when
both endpoints remain in the top $k$; confidence and partner conflicts are
resolved before emission. It follows the non-LLM entity-search idea of
FastGraphRAG and is not a reproduction of a full GraphRAG pipeline with
LLM-based community summarization or generation.
\textbf{EPGM} denotes one training-free structured retriever, reported with
its directed dependency-path strategy: it walks dependency edges from dense
seeds under schema gating, scores paths over schema-level semantic roles, and
promotes each accepted evidence partner to just after its anchor, so graph
structure reorders candidates without changing any dense score. The typed-beam
and query-conditioned diffusion strategies remain reproducibility
diagnostics only. Neither
the reported default nor the diagnostics use task-specific parameter
training. \textbf{EPGM R-GCN} is its learned
relation-aware counterpart. It applies three layers of relation-specific
message passing to frozen E5-base-v2 node representations, ranks
\texttt{ToolOutput} candidates, and jointly predicts logical
output-to-output dependency edges. Dense-FT is trained for one epoch under the reported protocol (not an exhaustive epoch sweep).
EPGM R-GCN is trained for up to 15 epochs and selects a checkpoint on the
complete 2,850-item development split with
$0.50\,\mathrm{FS@5}+0.25\,\mathrm{MRR}+0.25\,\mathrm{EdgeF1@10}$;
ties are resolved by FS@5, Edge F1@10, MRR, and then the earlier epoch.
Dense-FT remains a seed-13 snapshot, while the EPGM R-GCN full model and
the completed ablations (\emph{w/o graph}, \emph{w/o edge type},
\emph{w/o edge weight}, and \emph{w/o edge rerank}) are summarized over
seeds $\{13,17,29\}$ in Table~\ref{tab:provenance-three-seed}. Interim
seed-13 main/ablation tables remain for single-seed readability; the full
main-table row is identical to the full ablation row on seed 13. The real-trace
study is qualitative and uses training-free retrievers only; EPGM R-GCN is
not applied there.

\paragraph{Metrics.}
Recall@$k$ measures the fraction of gold evidence retrieved in the
top $k$, while Full Support@$k$ is one only when all gold evidence is
present, so it directly measures whether the complete support set is
recovered. MRR measures the rank of the first gold item. Connected
Evidence Recall is one only when all gold nodes enter the top $k$ and those
gold nodes are connected within the top-$k$ induced graph. Path Recall@10
is a per-query binary measure requiring both (i) all gold dependency
endpoints in the top 10 and (ii) every directed gold dependency reachable
in the returned subgraph. It is averaged over queries with a nonempty gold
dependency set. Edge Recall@10 measures direct recovery of gold
dependency edges. Edge metrics are micro-aggregated against all gold dependencies. Edge
Precision@10 divides true-positive edges by emitted edges whose endpoints
are both in the top 10; Edge Recall@10 divides by all gold edges, so a gold
edge with a missing top-10 endpoint remains a false negative. Edge F1@10
combines these two quantities. Abstention rate is the fraction of considered top-5 source outputs
that produce no accepted dependency after pool, threshold, and conflict
filtering. Average retrieved edges counts public logical output-to-output
edges, not the two native \texttt{feeds}/\texttt{returns} edges used to
reconstruct each trace. HotpotQA does not provide directed gold dependency
paths, so we report connectivity rather than path and edge metrics for that
dataset. A dash in the path columns indicates that a method does not return a
retrieved evidence subgraph with visible candidate-to-candidate edges, so
the metric is undefined rather than zero.
To further distinguish ``retrieving part of the evidence on average''
from ``recovering the complete support chain,'' we also use the
completeness gap
$\Delta_{\mathrm{comp}}@k=\mathrm{Recall}@k-\mathrm{Full\ Support}@k$
as a diagnostic quantity. A smaller gap means that ordinary recall is
less inflated by examples in which only part of the gold evidence is
retrieved. It is used to interpret the primary metrics rather than as
an additional headline metric.

\subsection{Large-Scale Evidence Retrieval (RQ1)}

\paragraph{Graph encoding yields higher complete-support point estimates
on all three datasets.}
Table~\ref{tab:hotpot} reports HotpotQA results. In the reported seed-13
run, Dense-FT R-GCN has the highest point estimate on every metric.
Relative to Dense-FT, relational graph encoding raises Full Support@5
from 68.96\% to 79.98\% and Full Support@10 from 86.75\% to 92.79\%.
Thus, even after task-specific dense fine-tuning, explicit evidence
relations are associated with more complete support sets under a limited
retrieval budget.

\begin{table*}[t]
\centering
\small
\setlength{\tabcolsep}{4pt}
\caption{HotpotQA evidence retrieval point estimates (\%, seed 13).}
\label{tab:hotpot}
\begin{tabular}{lrrrrrrr}
\hline
Method & R@5 & R@10 & FS@5 & FS@10 & MRR & Conn.@5 & Conn.@10 \\
\hline
BM25 & 67.88 & 81.27 & 37.50 & 60.04 & 82.70 & 26.32 & 45.65 \\
Dense & 76.98 & 88.86 & 53.20 & 75.38 & 85.30 & 37.73 & 59.44 \\
BM25-Graph & 71.94 & 86.10 & 44.90 & 70.11 & 82.22 & 35.20 & 57.21 \\
Dense-Graph & 77.77 & 89.46 & 54.97 & 77.06 & 85.24 & 40.65 & 62.21 \\
Dense R-GCN & 88.71 & 96.00 & 75.99 & 90.90 & 91.56 & 57.11 & 73.43 \\
Dense-FT & 85.48 & 94.36 & 68.96 & 86.75 & 89.32 & 49.69 & 69.14 \\
\textbf{Dense-FT R-GCN} & \textbf{90.88} & \textbf{96.87} &
\textbf{79.98} & \textbf{92.79} & \textbf{93.05} &
\textbf{59.82} & \textbf{75.24} \\
\hline
\end{tabular}
\end{table*}

The observed Full Support@5 difference is larger on 2WikiMultiHopQA
(Table~\ref{tab:2wiki}). Dense-FT R-GCN reaches 94.44\%, 13.01 points
above Dense-FT, and obtains 93.99\% Path Recall@10 and 88.92\% Edge
Recall@10. Compared with Dense-Graph, the path and edge point estimates
are higher by 35.86 and 31.70 points. Dense R-GCN also yields higher
point estimates than the legacy Dense-Graph reranker, suggesting that the
observed gain is not attributable solely to dense-encoder fine-tuning.

\begin{table*}[t]
\centering
\small
\setlength{\tabcolsep}{4pt}
\caption{2WikiMultiHopQA evidence-path retrieval point estimates (\%, seed 13).}
\label{tab:2wiki}
\begin{tabular}{lrrrrrrr}
\hline
Method & R@5 & R@10 & FS@5 & FS@10 & MRR & Path@10 & Edge@10 \\
\hline
BM25 & 52.33 & 67.74 & 18.46 & 36.24 & 65.93 & -- & -- \\
Dense & 67.76 & 79.53 & 36.28 & 55.86 & 86.33 & -- & -- \\
BM25-Graph & 54.98 & 73.45 & 22.71 & 47.22 & 62.37 & 38.41 & 38.20 \\
Dense-Graph & 68.07 & 81.53 & 37.61 & 60.46 & 87.06 & 58.13 & 57.22 \\
Dense R-GCN & 93.51 & 97.42 & 87.64 & 95.06 & 95.25 & 89.43 & 85.61 \\
Dense-FT & 91.90 & 98.14 & 81.43 & 96.15 & 94.44 & -- & -- \\
\textbf{Dense-FT R-GCN} & \textbf{97.06} & \textbf{99.06} &
\textbf{94.44} & \textbf{98.24} & \textbf{97.46} &
\textbf{93.99} & \textbf{88.92} \\
\hline
\end{tabular}
\end{table*}

MuSiQue-Ans exhibits the same pattern (Table~\ref{tab:musique}). In the
reported run, Dense-FT R-GCN improves Full Support@5 from 58.42\% to
62.75\% and reaches 80.75\% Path Recall@10, 13.35 points above
Dense-Graph. It also has the highest MRR point estimate (92.21\%), so the
observed benefit spans both first-hit ranking and complete evidence/path
recovery on this dataset.

\begin{table*}[t]
\centering
\small
\setlength{\tabcolsep}{4pt}
\caption{MuSiQue-Ans v1.0 evidence-path retrieval point estimates (\%, seed 13).}
\label{tab:musique}
\begin{tabular}{lrrrrrrr}
\hline
Method & R@5 & R@10 & FS@5 & FS@10 & MRR & Path@10 & Edge@10 \\
\hline
BM25 & 55.20 & 70.18 & 18.83 & 37.92 & 75.44 & -- & -- \\
Dense & 68.66 & 85.07 & 37.09 & 64.21 & 85.78 & -- & -- \\
BM25-Graph & 60.48 & 77.09 & 29.32 & 52.95 & 75.93 & 51.02 & 53.41 \\
Dense-Graph & 71.18 & 87.20 & 44.34 & 69.54 & 84.29 & 67.40 & 68.80 \\
Dense R-GCN & 80.23 & 93.63 & 56.91 & 83.46 & 88.23 & 79.34 & 78.18 \\
Dense-FT & 80.79 & 93.75 & 58.42 & 83.36 & 91.48 & -- & -- \\
\textbf{Dense-FT R-GCN} & \textbf{83.89} & \textbf{94.60} &
\textbf{62.75} & \textbf{85.45} & \textbf{92.21} &
\textbf{80.75} & \textbf{78.96} \\
\hline
\end{tabular}
\end{table*}

Across all three datasets, graph encoding primarily improves complete
support and path recovery. In particular, Dense-FT and Dense-FT
R-GCN differ by less than one point on 2Wiki Recall@10, but by more
than thirteen points on Full Support@5. This pattern indicates that the
model does more than raise average recall of isolated relevant sentences:
it more often brings all members of a reasoning chain into the top-$k$
context. Heuristic graph reranking sometimes helps, but its improvements
are less uniform across datasets and metrics than those of the learned
relation-aware models in these runs.

Because the evidence-graph construction rules differ across datasets, we
use RQ1 to evaluate the overall structured-retrieval hypothesis and reserve
component attribution for the controlled execution-provenance ablations in
RQ2.

\subsection{Controlled Execution-Provenance Retrieval (RQ2)}

\paragraph{Execution structure primarily improves dependency-set
completion.}
Table~\ref{tab:provenance} compares sparse retrieval, dense retrieval,
lightweight entity-graph propagation, training-free provenance search,
and learned provenance encoding. Without task-specific training, the
EPGM dependency-path preset raises Full Support@5 and Full Support@10
over Dense by 29.39 and 21.29 points and Recall@5 by 13.24 points, while
MRR declines slightly (down 0.70 points). Relative to
FastGraphRAG-style, this preset improves Full Support@5 from 32.90\% to
54.09\% and Full Support@10 from 56.12\% to 73.48\%; the three methods
tie on Recall@2 (all 46.12\%), and MRR differs by less than one point
(80.90\% versus 80.21\%). The gains are consistent rather than
averaging artifacts: per query, dependency-path search improves
Recall@5 on 994 queries and degrades it on 237, and improves Full
Support@5 on 984 against 146.

Two properties of this comparison are worth stating plainly. First, the
method permutes the entire candidate list rather than reordering the dense
top-10, which is why it also lifts Recall@10 (by
10.56 points over Dense) instead of only redistributing rank mass inside
a fixed set. Second, the improvement is concentrated in set-completion
metrics and does not extend to first-hit quality: Recall@2 is unchanged
and MRR moves in the opposite direction, degrading on 236 queries against
39 improved. The mechanism recovers dependency partners that lie outside
the dense top-5, at the cost of occasionally displacing an output that
dense retrieval had already placed first.

\begin{table*}[t]
\centering
\small
\resizebox{\textwidth}{!}{%
\begin{tabular}{lrrrrrrrrrrr}
\hline
Method & R@2 & R@5 & R@10 & F1@5 & F1@10 & FS@5 & FS@10 & MRR & Path@10 & Edge@10 & ms/query \\
\hline
BM25 & 28.32 & 44.05 & 58.59 & 25.17 & 19.53 & 7.33 & 23.68 & 55.16 & -- & -- & \textbf{0.60} \\
Dense & 46.12 & 61.54 & 76.04 & 35.17 & 25.35 & 24.69 & 52.19 & 80.91 & 0.00 & 0.00 & 37.51 \\
FastGraphRAG-style (non-LLM) & 46.12 & 65.59 & 78.01 & 37.48 & 26.00 & 32.90 & 56.12 & 80.90 & 10.24 & 10.10 & 105.53 \\
EPGM dependency path (training-free) & 46.12 & 74.78 & 86.60 & 42.73 & 28.87 & 54.09 & 73.48 & 80.21 &
44.93 & 44.12 & 39.04 \\
\hline
Dense-FT & \textbf{71.03} & 91.32 & 97.91 & 52.18 & 32.64 & 83.02 & 95.83 & 93.09 & -- & -- & 19.53 \\
\textbf{EPGM R-GCN} & 57.86 & \textbf{98.09} & \textbf{99.58} &
\textbf{56.05} & \textbf{33.19} & \textbf{96.74} & \textbf{99.30} & \textbf{95.09} &
\textbf{90.85} & \textbf{90.85} & 68.69 \\
\hline
\end{tabular}%
}
\caption{Interim seed-13 synthetic 2Wiki execution-provenance retrieval
results. Quality metrics are percentages; latency is milliseconds per
query. The trainable rows will be updated after the three-seed matrix is
complete.}
\label{tab:provenance}
\end{table*}

In the current seed-13 snapshot, task adaptation is itself a strong
baseline: Dense-FT reaches 91.32\% Recall@5 and 83.02\% Full Support@5.
EPGM R-GCN improves these values to 98.09\% and 96.74\%, gains of 6.77
and 13.72 points;
it also raises Full Support@10 by 3.47 points and MRR by 2.00 points.
Unlike the flat Dense-FT retriever, the learned provenance model emits a
dependency subgraph and recovers 90.85\% of gold paths and 90.85\% of
gold dependency edges at $k=10$, with Edge Precision@10 / Edge F1@10 of
84.70\% / 87.66\% under the sparse decoder. Its Recall@2 (57.86\%) is
below Dense-FT (71.03\%): the structured model is not primarily an
early-rank specialist, but a complete-support and dependency decoder.
Table~\ref{tab:provenance-ablation} ablations use the same protocol and
splits as this full row, so within-table differences localize components
of the reported EPGM R-GCN rather than a second training regime. Path
Recall and Edge Recall remain recall measures and do not by themselves
show whether the predicted edge set is sparse or well calibrated; Edge
Precision and Edge F1 supply that check.

\paragraph{Structural modeling substantially narrows the gap between
partial and complete retrieval.}
Because every synthetic trajectory contains two gold tool outputs, Full
Support can be converted directly into the number of completely recovered
trajectories. Table~\ref{tab:provenance-budget} reports the complete counts
at top-5 and top-10, the additional trajectories completed when the budget
is increased, the number still incomplete at top-10, and
$\Delta_{\mathrm{comp}}@5$. Rather than duplicating the main table, this
view maps the percentage metrics back to the 2,851 test queries and shows
how many complete dependency chains each method actually resolves.

\begin{table*}[t]
\centering
\small
\setlength{\tabcolsep}{5pt}
\caption{Interim seed-13 complete-support and retrieval-budget diagnostics
on synthetic execution trajectories. Complete-count columns use 2,851 as
their denominator; the completeness gap is measured in percentage points.}
\label{tab:provenance-budget}
\resizebox{\textwidth}{!}{%
\begin{tabular}{lrrrrr}
\hline
Method & Complete FS@5 & Complete FS@10 & Added at $k:5\rightarrow10$ & Incomplete FS@10 & $\Delta_{\mathrm{comp}}@5$ \\
\hline
BM25 & 209 & 675 & 466 & 2,176 & 36.72 \\
Dense & 704 & 1,488 & 784 & 1,363 & 27.50 \\
FastGraphRAG-style (non-LLM) & 938 & 1,600 & 662 & 1,251 & 23.22 \\
EPGM dependency path (training-free) & 1,542 & 2,095 & 553 & 756 & 19.40 \\
\hline
Dense-FT & 2,367 & 2,732 & 365 & 119 & 8.30 \\
\textbf{EPGM R-GCN} & \textbf{2,758} & \textbf{2,831} & 73 &
\textbf{20} & \textbf{1.35} \\
\hline
\end{tabular}%
}
\end{table*}

Among methods without task-specific training, the EPGM dependency-path
preset completely recovers 1,542 trajectories at top-5, 838 more than
Dense and 604 more than FastGraphRAG-style; its completeness gap also
decreases from 27.50 and 23.22 points to 19.40 points. In this controlled
benchmark, execution-dependency search is therefore better at bringing
both relevant outputs into a limited context, although
756 trajectories remain incomplete at top-10 and the training-free preset
does not solve the task. Dense completes 784 additional trajectories when
the budget grows from 5 to 10, the largest increase among all methods, but
starts from the lowest small-budget count; the EPGM dependency-path preset
adds the fewest (553) precisely because it has already completed most of
what a larger budget would recover, indicating that it relies on
completing dependencies at small $k$ rather than on expanding the
retrieval budget.

Among trained methods, EPGM R-GCN completely recovers 391 more trajectories
than Dense-FT at top-5 and 99 more at top-10, reducing the residual top-10
failures from 119 to 20. More importantly,
$\Delta_{\mathrm{comp}}@5$ shrinks from 8.30 to 1.35 points, while
$\Delta_{\mathrm{comp}}@10$ shrinks from 2.09 to 0.28 points. The main
value of relational propagation is therefore not only higher average gold
output recall, but converting partial hits that recover one output into
complete dependency sets that provide a more complete input for downstream
reasoning.

We next ablate the same EPGM R-GCN protocol to localize which structural
components drive these gains.

\paragraph{Model and ablation protocol.}
We use the EPGM R-GCN from Section~\ref{sec:evidence-tracing} on the shared
74,012/2,850/2,851 splits (seeds 13, 17, and 29; joint checkpoint rule;
sparse decoder). The detailed rows currently populated below are the
completed seed-13 snapshot; seeds 17 and 29 will be aggregated only after
all full/ablation pairs pass the same artifact-identity checks.
The full model uses 714,245 pairs. We
compare the full model to three training ablations and reserve one paired
inference-time ablation.

\emph{w/o graph} sets the number of graph-convolution layers to zero while
retaining the input projection, node-type embedding, candidate and edge
heads, logical-transition enumeration, and structured decoder.
\emph{w/o edge type} retains topology, relation IDs, and weights but shares
one linear message transform across all relation IDs. \emph{w/o edge
weight} changes only \texttt{feeds} messages: each source's two artifact
weights are replaced by their mean (0.75/0.75 under this construction), so
endpoints, relation IDs, directions, and total source message mass are
preserved; non-\texttt{feeds} weights are unchanged. \emph{w/o edge rerank} must reuse the exact full-model
pairs and checkpoint for the same seed, disable only target promotion, and
retain thresholded edge diagnostics. It will enter the final table only when
that checkpoint identity is verified; the earlier unpaired seed-13 artifact
is not used for causal attribution.

Within each seed, every training ablation selects its own development
checkpoint under the joint objective. The current full seed-13 row matches
Table~\ref{tab:provenance}. For the final analysis, each variant's metrics
will be summarized as mean $\pm$ sample standard deviation across the three
seeds. For every ablation and metric, predictions will be paired with the
full model by seed and task ID; full-minus-ablation per-query deltas will be
pooled across matched seed-query pairs and summarized with a 95\% percentile
bootstrap interval (2,000 resamples; bootstrap seed 13). Positive deltas
therefore favor the full model. Full-Support@5/@10 discordant counts
(full-only versus ablation-only successes) will be reported alongside the
intervals in the supplementary material.

\paragraph{Interim seed-13 metrics show broad gains from graph propagation
across objectives.}
Table~\ref{tab:provenance-ablation} reports six metrics spanning early
ranking, complete support, dependency recall, and dependency quality:
Recall@2 measures early ranking under a very small budget, Full Support@5
tests whether both relevant outputs enter the context together, Path@10
and Edge@10 measure dependency-structure recall, and Edge Precision@10 and
Edge F1@10 measure the quality of predicted dependency edges rather than
recall alone. Under this sparse, precision-oriented decoder, the full model is the best
variant on Full Support@5, Edge Precision@10, and Edge F1@10, reaching
96.74\%, \textbf{84.70\%}, and \textbf{87.66\%}, respectively. Removing
graph propagation reduces these three to 84.74\%, 53.84\%, and 64.34\%,
while Full Support@5 and MRR drop by 12.00 and 12.09 points, showing that
the structural gain is not confined to one evaluation criterion. Recall@2 is not the main beneficiary of graph
propagation (the full model and \emph{w/o graph} differ by only 0.83
points); dependency-edge precision and complete support separate them.
This pattern matches the main-table contrast with Dense-FT, where EPGM
R-GCN trails on Recall@2 but leads on Full Support and dependency-edge
quality.

\begin{table*}[t]
\centering
\small
\setlength{\tabcolsep}{3.5pt}
\resizebox{\textwidth}{!}{%
\begin{tabular}{lrrrrrrrrrr}
\hline
Variant & R@2 & R@5 & R@10 & FS@5 & FS@10 & MRR & Path@10 & Edge@10 & EdgeP@10 & EdgeF1@10 \\
\hline
Full EPGM R-GCN & 57.86 & \textbf{98.09} & 99.58 & \textbf{96.74} & 99.30 & 95.09 &
90.85 & 90.85 & \textbf{84.70} & \textbf{87.66} \\
w/o graph & 57.03 & 91.07 & 96.97 & 84.74 & 94.11 & 83.00 & 80.01 & 79.94 &
53.84 & 64.34 \\
w/o edge type & 68.84 & 97.04 & 99.37 & 95.12 & 98.91 & 91.86 & 89.23 & 89.20 &
74.86 & 81.40 \\
w/o edge weight & 60.01 & 97.84 & \textbf{99.70} & 96.53 & \textbf{99.51} &
94.55 & 88.50 & 88.46 & 80.34 & 84.21 \\
\hline
\end{tabular}%
}
\caption{Interim seed-13 EPGM R-GCN component ablations on the synthetic
2Wiki execution-provenance benchmark (\%). EdgeP@10 and EdgeF1@10 are
Edge Precision@10 and Edge F1@10. Rows share the main-table full-model
protocol; final claims will use Table~\ref{tab:provenance-three-seed}.}
\label{tab:provenance-ablation}
\end{table*}

In the seed-13 snapshot, removing graph message passing causes the largest
observed quality degradation. Relative to the full model, \emph{w/o graph} drops
Full Support@5, MRR, and Path Recall@10 by 12.00, 12.09, and 10.84 points,
and Edge Precision@10 and Edge F1@10 by 30.86 and 23.32 points. In sample
counts, \emph{w/o graph} completely recovers 342 fewer trajectories at
top-5 and leaves 148 more trajectories incomplete at top-10; yet its
Recall@2 is nearly tied with the full model, further showing that under
sparse decoding the gain from graph propagation concentrates in complete
support and dependency-edge quality rather than the first hit under a very
small budget. Because this variant retains the same frozen text
representations, node-type embeddings, candidate MLP, and dependency-edge
MLP, the controlled difference attributes the largest observed gain to
neighborhood information exchange rather than to the text encoder or output
heads. Three-layer propagation allows upstream evidence and the query anchor
to influence downstream outputs through the call hierarchy and
\texttt{feeds}/\texttt{returns} connections, which is consistent with the
improved joint recovery of both gold outputs under a small budget.

Replacing relation-specific transforms with a shared transform lowers the
quality of predicted dependency edges: Edge Precision@10 and Edge F1@10
fall by 9.84 and 6.26 points, MRR by 3.23 points, and Full Support@5 by
1.62 points (46 fewer complete top-5 trajectories, and incomplete top-10
trajectories rise from 20 to 31). Notably, this variant's Recall@2 is
10.98 points \emph{above} the full model, while its dependency-edge
precision is lower. Untyped propagation therefore favors the very early
candidate ranking on this split but does not preserve the full model's
edge-selection quality. This variant retains exactly the same
adjacency, propagation depth, and edge weights; the only difference is
whether $W_r$ distinguishes relations. The observed single-seed difference
is consistent with a contribution from typed propagation to dependency-edge
quality: the directions and semantics carried by call-hierarchy edges and
field-level data-flow edges are not fully replaced by one untyped
transformation, particularly when
separating structurally adjacent but functionally different \texttt{Agent},
\texttt{ToolCall}, and \texttt{ToolOutput} messages.

The evidence for edge weights also concentrates on precision and the path
side. Replacing each source's two \texttt{feeds} confidences by their mean
lowers Edge Precision@10, Edge F1@10, and Path Recall@10 by 4.36, 3.45,
and 2.35 points and MRR by 0.54 points; at the same time, Full Support@5
dips by only 0.21 points while Full Support@10 rises by 0.21 points (6 fewer
and 6 more complete examples among 2,851 test queries, respectively).
Calibrated within-source weights therefore mainly improve dependency-edge
quality, path recovery, and MRR, while their effect on top-$k$ complete
support is approximately neutral on this split. This
matches the implementation: weights directly scale propagated messages and
can alter fine-grained candidate order and dependency-edge logits, but
Full Support does not change as long as both gold outputs remain inside the
same top-$k$ set.

\paragraph{Development selection, residual failures, and output density
provide consistent structural evidence.}
Table~\ref{tab:provenance-ablation-diagnostics} decomposes test
percentages into model-selection behavior, complete-recovery counts, and
output-subgraph density. Dev FS@5 is the Full Support@5 value at the
checkpoint selected by the joint objective on all 2,850 development
queries; it is not an independently maximized FS@5. Average edge count is
measured when every model returns ten nodes and counts logical
output-to-output edges. The decoder remains sparse across variants, with
1.07--1.48 emitted edges per query. Because the evaluator stores
at most 50 detailed failures, the incomplete counts are computed directly
from the complete test-set metrics rather than estimated from the failure
file.

\begin{table*}[t]
\centering
\small
\setlength{\tabcolsep}{5pt}
\caption{Interim seed-13 model selection, complete coverage, and
output-density diagnostics for the EPGM R-GCN ablations. Dev FS@5 is a
percentage; the other complete and incomplete columns use 2,851 as their
denominator.}
\label{tab:provenance-ablation-diagnostics}
\resizebox{\textwidth}{!}{%
\begin{tabular}{lrrrrrr}
\hline
Variant & Best epoch & Dev FS@5 & Complete FS@5 & Complete FS@10 & Incomplete FS@10 & Avg. edges \\
\hline
Full EPGM R-GCN & 14 & \textbf{96.88} & \textbf{2,758} & 2,831 & 20 & 1.07 \\
w/o graph & 15 & 85.16 & 2,416 & 2,683 & 168 & 1.48 \\
w/o edge type & 14 & 94.42 & 2,712 & 2,820 & 31 & 1.19 \\
w/o edge weight & 13 & 96.32 & 2,752 & \textbf{2,837} & \textbf{14} & 1.10 \\
\hline
\end{tabular}%
}
\end{table*}

Development and test results provide evidence in the same direction. The
selected-checkpoint Dev FS@5 of \emph{w/o graph} is 11.72 points below that
of the full model, and that of \emph{w/o edge type} is 2.46 points lower; on the test set these
variants leave 168 and 31 top-10 trajectories incomplete, respectively,
compared with only 20 for the full model. The full and three training
ablations select checkpoints by the development
joint score, with selected epochs of 14 (full), 15 (\emph{w/o graph}), 14
(\emph{w/o edge type}), and 13 (\emph{w/o edge weight}); their corresponding
Dev FS@5 values preserve the same
ordering on the test set. Source-mean edge weighting trails
the full model on Dev FS@5 by only 0.56 points, while its small
test-set reversals in local metrics (such as \emph{w/o edge weight} on
FS@10) do not appear on development. This mismatch further
suggests that differences of only a few tenths of a point more likely
reflect split or seed variation than a stable advantage. However, most
variants attain their best checkpoints near the 15-epoch training limit, so
whether longer training would close some gaps requires additional
experiments; the results compare components under a common training
budget but do not establish that every variant has fully converged.

\subsubsection{Three-Seed Robustness Report}

Table~\ref{tab:provenance-three-seed} reports the completed three-seed matrix
for the full model and three training ablations plus one paired
inference-time ablation, computed with
\texttt{scripts/analyze\_provenance\_ablation.py} (2,000 bootstrap resamples;
bootstrap seed 13). A confidence interval that crosses zero is treated as
inconclusive rather than as evidence that the full model is better.

\begin{table*}[t]
\centering
\small
\setlength{\tabcolsep}{2.6pt}
\resizebox{\textwidth}{!}{%
\begin{tabular}{lrrrrrrr}
\hline
Variant & FS@5 $\mu{\pm}\sigma$ & MRR $\mu{\pm}\sigma$ & Path@10 $\mu{\pm}\sigma$ & EdgeP@10 $\mu{\pm}\sigma$ & EdgeF1@10 $\mu{\pm}\sigma$ & $\Delta$FS@5 [95\% CI] & $\Delta$EdgeF1 \\
\hline
Full EPGM R-GCN & $97.08\pm0.21$ & $95.13\pm0.10$ & $91.32\pm0.39$ & $84.19\pm0.65$ & $87.61\pm0.41$ & -- & -- \\
w/o graph & $84.51\pm0.60$ & $83.33\pm0.91$ & $78.53\pm0.39$ & $56.78\pm1.33$ & $65.88\pm0.86$ & $12.57$ $[11.77, 13.32]$ & $21.73$ $[21.44, 22.29]$ \\
w/o edge type & $95.12\pm1.13$ & $92.60\pm0.45$ & $88.85\pm2.17$ & $75.75\pm3.08$ & $81.69\pm1.12$ & $1.95$ $[1.54, 2.37]$ & $5.92$ $[5.41, 6.77]$ \\
w/o edge weight & $96.48\pm0.28$ & $94.97\pm0.49$ & $89.54\pm1.05$ & $80.95\pm0.30$ & $85.00\pm0.37$ & $0.60$ $[0.23, 0.96]$ & $2.61$ $[2.37, 2.81]$ \\
w/o edge rerank & $96.64\pm0.17$ & $95.09\pm0.11$ & $91.32\pm0.39$ & $84.21\pm0.63$ & $87.62\pm0.40$ & $0.43$ $[0.27, 0.60]$ & $-0.01$ $[-0.01, 0.00]$ \\
\hline
\end{tabular}%
}
\caption{Three-seed robustness report over seeds 13, 17, and 29 (\%). Metric
cells are mean $\pm$ sample standard deviation. $\Delta$FS@5 is the mean
full-minus-ablation percentage-point difference with a query-paired 95\%
percentile-bootstrap interval over matched seed-query pairs ($n{=}8553$).
$\Delta$EdgeF1 is the mean seed-level micro-average Edge F1@10 difference
with the min--max range across the three seeds, because Edge F1@10 is
micro-aggregated rather than stored per query. Positive deltas favor the
full model.}
\label{tab:provenance-three-seed}
\end{table*}

\paragraph{Residual errors point to second-hop coverage rather than
indiscriminate subgraph expansion.}
Manual inspection of all 20 residual top-10 failures shows that in 16
cases the model recovers one of the two gold outputs and misses the other;
only four cases miss both. The remaining errors are therefore mostly
``near-complete but missing one hop,'' and future improvements should
prioritize seed coverage and joint ranking for the second dependency
output rather than merely expanding the returned subgraph. Because a
dependency edge is decoded only after both endpoints enter the top-10
set, these 16 single-endpoint misses also explain why small changes in
candidate recall are simultaneously reflected in Path@10 and Edge@10.

Output density further argues against, within these ablations, the
alternative explanation that predicting more edges is sufficient to raise
edge recall. \emph{w/o graph}
returns 1.48 edges on average, the most of any variant, yet is the lowest
on both Edge Recall@10 and Edge Precision@10 (79.94\% and 53.84\%); \emph{w/o
edge weight} returns 1.10 edges on
average, also more than the full model's 1.07, yet achieves lower Edge
Recall@10. The full model therefore obtains the highest Edge Precision@10 and Edge
F1@10 while emitting the fewest average edges among the training ablations.
This pattern indicates that typed transforms and calibrated weights improve
dependency selection rather than merely making the decoder connect more
candidates.
The decoder keeps at most one
successor per source, which matches the two-hop chain structure of this
synthetic benchmark; natural trajectories in which one output is consumed
by multiple calls will additionally require evaluation of branching
decoders and edge precision.

Taken together, the three-seed matrix in Table~\ref{tab:provenance-three-seed}
supports the following hierarchy. Graph message passing remains the dominant
component: \emph{w/o graph} drops mean Full Support@5 by 12.57 points with
query-paired 95\% CI $[11.77, 13.32]$ and lowers Edge F1@10 by 21.73 points
across seeds. Relation-specific transforms are second: \emph{w/o edge type}
drops Full Support@5 by 1.95 $[1.54, 2.37]$ and Edge F1@10 by 5.92 points.
Source-calibrated edge weights contribute a smaller but consistently positive
Full Support@5 delta of 0.60 $[0.23, 0.96]$ and a 2.61-point Edge F1@10 gap.
Disabling target promotion (\emph{w/o edge rerank}) leaves Path@10 and Edge
F1@10 essentially unchanged while yielding only a 0.43-point Full Support@5
gain for the full model.

\paragraph{Controlled mechanism example.}
To illustrate how structured provenance converts a partial hit into a
complete hit under the RQ2 protocol, we analyze a reasoning-type
provenance query derived from 2WikiMultiHopQA: ``Who is Anne Woodville's father-in-law?'' (gold
answer: Edmund Grey, 1st Earl of Kent). The query requires a two-hop
dependency: first Anne Woodville's spouse George Grey, 2nd Earl of Kent
is obtained through the \texttt{spouse} relation (first-hop output), and
then that spouse's father Edmund Grey is obtained through the
\texttt{father} relation (second-hop output). In the execution-provenance
graph, each step corresponds to a \texttt{ToolOutput} node, and the
second-hop output depends on the first-hop output through
\texttt{feeds}/\texttt{returns} edges.

Dense retrieval hits the first-hop output within the top-10 but misses the
second-hop output, so it falls into a \texttt{missing\_full\_support}
failure: it finds the spouse information that is semantically related to
the query directly, but cannot locate the indirect ``spouse's father''
dependency by text similarity alone. EPGM R-GCN recovers the complete
dependency pair on the same query: three-layer relational propagation lets
the first-hop output representation flow to the second-hop candidate along
\texttt{feeds}/\texttt{returns} edges, the dependency head assigns this
logical transition a high score, and both gold outputs are brought into the
top-$k$ set while the \texttt{spouse}$\rightarrow$\texttt{father}
dependency path is reconstructed. This is consistent with the aggregate
conclusion from Table~\ref{tab:provenance-budget}: the structural gain
concentrates on completing the second hop of the dependency set rather
than on improving the rank of the first relevant output.


\subsection{Real Multi-Agent Trace Study (RQ3)}

\paragraph{Setting: one naturalistic session, qualitative comparison.}
Automatic ranking metrics are primary for RQ1--RQ2. RQ3 instead inspects
how retrievers behave on SupplyChainTrace, a single web-enabled
supply-chain security session over seven pinned PyPI dependencies
(\texttt{requests==2.19.1}, \texttt{urllib3==1.24.1},
\texttt{PyYAML==3.13}, \texttt{Jinja2==2.10}, \texttt{Django==2.2.0},
\texttt{cryptography==2.3}, \texttt{Werkzeug==0.14.1}).
The session comprises a root coordinator, seven per-package subagents,
an independent verifier, and two verification workers. The recorded
provenance graph contains 150 retrievable nodes (task, agents, tool outputs,
observations, claims, verifications, decisions, and the final answer)
and includes genuine revision structure: early claims (Django
minimum-safe upgrade 4.2.26; a PyYAML-first ranking) are overturned via
\texttt{contradicts}/\texttt{invalidates} after verification, while
invalidated claims remain for audit.

We compare BM25, dense retrieval, FastGraphRAG-style entity search, and
the EPGM dependency-path preset on the same candidate pool, without
task-specific training or trace-specific tuning. Because a single session
cannot support stable leaderboard claims, we do \emph{not} treat aggregate
Recall or Full-Support tables as primary evidence for RQ3. On this trace the
preset emits typed provenance edges on 16 of 20 queries (56 edges, 38
\texttt{depends\_on} and 18 \texttt{supports}) against 5 of 20 single edges
for FastGraphRAG-style entity bridges and none for Dense, and it matches Dense
on Full Support@5 (4/20) and Full Support@10 (7/20). Its ranking metrics are
lower than Dense: Recall@5 47.83 against 50.50 and MRR 57.92 against 65.83.
That gap should not be over-read, but neither should it be hidden. Of the 56
promoted partners only 13 are gold, so promotion is 23.2\% precise here, and
the recorded \texttt{depends\_on} graph is dense enough that a bounded walk
reaches many non-evidential neighbours. At the same time only 3 of 20 queries
change their top-5 gold membership at all (one better, two worse), so the
aggregate difference rests on two queries and is not separable from noise at
this sample size. We therefore report no ranking benefit on this trace and
rest the RQ3 claim on what the typed edges expose, which the analysis below
illustrates.

\paragraph{Revision query: typed edges surface overturned claims for
audit.}
Consider a query that asks what Django's minimum-safe version was first
judged to be, why that judgment changed, and what the final decision is.
Dense retrieval already places the verification and the invalidated claim
\texttt{claim\_537BCA895F} (initially asserting 4.2.26) inside the top 10,
but returns them as independent records: nothing in the output states that
one overturns the other. The dependency-path preset returns the typed edge
\texttt{verify\_28A134BF07}$\xrightarrow{\texttt{contradicts}}$
\texttt{claim\_537BCA895F} and exposes
\texttt{lifecycle\_state=invalidated}, so the revision relation is
recoverable from the retrieval output itself rather than reconstructed by
the reader. FastGraphRAG-style entity bridges connect the two nodes by
shared surface entities but do not encode revision semantics. The point is
mechanistic and about interpretability: provenance structure answers an
audit question in the form the user asked, not that a 20-query average
score is decisive.

\paragraph{Support-chain query: typed edges expose the support set as a
structure.}
Query Q10 asks why version-range checks must use PEP~440 / SpecifierSet
semantics rather than naive string comparison. Three gold support nodes
are required, two of which (\texttt{claim\_0144C746E7} and
\texttt{obs\_42BF2FFDED}) have little surface overlap with the query.
FastGraphRAG-style returns only one of the three in its top 10. The
dependency-path preset emits the
$\texttt{obs}\xrightarrow{\texttt{supports}}\texttt{claim}$ edges that tie
support nodes to the claim they underwrite, so a consumer can tell which
retrieved records form one support set and which are unrelated hits. The same
pattern appears on a related cryptography upgrade query. On this trace, then,
typed traversal contributes an explicit evidence structure over the retrieved
set; the corresponding ranking gains appear only in the controlled RQ2
setting, where the candidate pool is large, dependency partners routinely sit
below the top-5 cutoff, and recorded \texttt{feeds} confidences give the path
score real dynamic range.

\paragraph{What this study does and does not show.}
SupplyChainTrace shows that the full EPGM schema can be populated from a
real multi-agent run with tool calls, parallel handoffs, and evidence
revision, and that typed retrieval can surface support and contradiction
edges that the compared rankings leave implicit. It does not show a ranking
improvement on this trace: the preset scores below Dense on Recall@5 and MRR,
and its promoted partners are only 23.2\% precise, so its contribution here is
the interpretable edge set rather than a better ordering. It does not provide
a scored real-world benchmark, does not evaluate EPGM R-GCN transfer, and does
not replace RQ1--RQ2 automatic results. Expanding to more sessions and
end-to-end tasks remains future work.

\subsection{Cross-Setting Analysis}

Across the three protocols, relational structure helps most when
relevance is split across dependent artifacts rather than concentrated in
a single surface-similar record. RQ1 yields higher complete-support point
estimates for relation-aware encoding on large multi-hop sets; RQ2 shows
the largest single-seed ablation drop when message passing is removed; and
RQ3's cases show the same schema supporting audit and support-path retrieval
on a real session where flat rankings leave relations implicit, even though
there the structure surfaces as typed edges rather than as a reordering.
The settings
are not interchangeable: multi-hop QA lacks native tool executions,
synthetic trajectories lack natural noise and true revision events, and
the real trace is too small for training or multi-seed claims. Read
together, they support applicability of the structured-retrieval
hypothesis under complementary constraints, not a universal real-world
SOTA result.

\subsection{Efficiency and Cost Analysis}

\paragraph{Structured retrieval increases observed latency while improving
complete-support recovery.}
The ms/query column of Table~\ref{tab:provenance} characterizes the
retrieval cost of each method. The reported training-free dependency-path preset averages 39.04 ms per
query, compared with 37.51 ms for Dense and 105.53 ms for the evaluated
FastGraphRAG-style implementation in this harness. Thus, bounded
execution-path search adds little observed latency beyond its dense base in
this run. Learned relation-aware propagation raises the EPGM R-GCN mean
latency to 68.69 ms per query on the delivered full run (median
67.61 ms over the 2,851 test queries), trading additional computation for
large gains in complete support and dependency recovery over Dense-FT. BM25
remains fastest at 0.60 ms per query, but its complete-support recovery is
far below all dense and graph methods. Baseline latencies were measured in
the same evaluation harness on the shared 2,851-query test split; the EPGM
R-GCN latency is the mean per-query time from the delivered full artifact.
These values are observed retrieval times rather than a controlled systems
benchmark with fixed hardware and caching, and they exclude data
generation, graph construction, and model training.

\subsection{Limitations and Robustness}

The experiments already provide partial evidence for structured
execution-provenance memory from three perspectives: controlled retrieval,
complete evidence recovery, and dependency reconstruction. Their scope is
still concentrated on offline retrieval over fully materialized static
graphs, so they provide only partial support for the full multi-agent
memory lifecycle. In particular, we do not measure streaming online graph
updates, mid-trajectory retrieval, or incremental index maintenance;
\textsc{UpdateGraph} is a schema-level description of event-local
construction, not an evaluated systems component. First, the synthetic
provenance benchmark is deterministically derived from labeled
question-answering evidence chains rather than collected from
heterogeneous, naturally occurring agent logs. The conversion also uses
supporting-fact information to materialize the required dependency as an
ordinary semantic head or rank-matched branch and rejects examples for
which a matched non-gold branch cannot be formed. The near-ceiling EPGM R-GCN numbers show recovery on a controlled,
label-derived topology where required \texttt{feeds} edges are inserted as
ordinary heads or rank-matched branches. This does \emph{not} by itself rule
out converter-specific shortcuts; controls such as shuffling \texttt{feeds}
endpoints at matched degree, topology-only baselines, or relation holdouts
are not yet reported. SupplyChainTrace is analyzed qualitatively for schema applicability and
case-level retrieval behavior, not as a scored multi-session benchmark.

Second, Dense-FT on the provenance benchmark remains a seed-13 snapshot; the
full
model and the other ablations are now summarized with three-seed
mean$\pm$std and paired Full Support@5 intervals in
Table~\ref{tab:provenance-three-seed}. Small single-seed reversals after
source-mean edge weighting should be read against those intervals rather
than as stable advantages. Removing graph propagation is the largest and
most stable cross-metric drop under the completed three-seed protocol.
The training-free dependency-path preset carries two tradeoffs that its
aggregate gains should not obscure. Its set-completion improvements do not
transfer to first-hit quality: Recall@2 is identical to Dense and MRR is
0.70 points lower, degrading on 236 of 2,851 queries against 39 improved,
because promoting a dependency partner into a small budget can displace an
output that dense retrieval had already ranked first. Its edge output also
trades precision for coverage, reaching 44.93\% Path Recall@10 against
10.24\% for entity-graph propagation but 12.34\% Edge Precision@10 against
34.33\%, since it emits 3.58 edges per query against 0.29. Applications
that consume the edge set directly would need a sparser decoder or a
confidence cutoff.

Furthermore, the same frozen configuration behaves very differently across
the two provenance settings, for a reason the schema makes explicit. Scoring
and gating are defined over semantic roles, but the two graphs populate those
roles with disjoint vocabularies and very different weight distributions. The
controlled benchmark records evidence hand-off only as \texttt{feeds} and
supplies calibrated confidences (over a thousand distinct values), so path
scores have real dynamic range and reorder results substantially. The recorded
trace uses \texttt{depends\_on}/\texttt{supports} and stores a constant
placeholder weight of 1.0 on every edge, so path scores collapse onto the
frozen relation priors and carry only three distinct levels. Role abstraction
is what makes one configuration \emph{applicable} to both; it does not make
the available signal equally informative. We report both rather than tuning
per dataset, but this means the retrieval gain is demonstrated on
label-derived topology with calibrated weights and the interpretability
benefit on the real trace, not both on both.
Third, the current experiments
evaluate retrieval, complete support, and path recovery, so they can
partially support the claim that structured memory provides more complete
grounding for downstream reasoning; system-level effects such as ``reducing
unsupported claims'' or ``improving downstream impact-tracing accuracy''
receive only indirect support until the retrieved subgraphs are connected
to real decision or generation tasks. The reported tables already include edge
precision and F1, showing that the structural components' gains lie in
dependency-edge quality rather than recall alone; adding branch-level
completeness would more fully quantify overprediction and multi-branch
recovery.

Fourth, \emph{w/o graph} removes only graph message passing while
retaining node-type embeddings, the candidate scorer, logical transitions
enumerated from graph structure, and the dependency-edge head, so it
answers whether propagation is necessary rather than capturing every
difference between a graph model and a fully flat model. The independently
trained Dense-FT in the main experiment already provides the flat text
baseline; to further localize each piece of structural information, future
work should separately remove node types, field-binding features, and the
edge-prediction loss, shuffle the data-flow relations, and evaluate
invalidation penalties on trajectories containing actual evidence-revision
events. Overall, the current results provide quantitative evidence for
compact support recovery and dependency decoding on the evaluated evidence
graphs and controlled provenance benchmark, plus qualitative applicability
evidence on
one naturalistic trace. They provide only partial support for a complete
multi-agent memory system; broader system-level conclusions require
multi-seed repetition, more natural trajectories, and end-to-end tasks.
